%lesson 2, 05/03/2026
\newpage
\subsection{Autovettori e autovalori}

Dato un operatore autoaggiunto $ A $, si consideri ora l'equazione agli autovalori
$$ A\ket{a} = a\ket{a} $$ 

Prima di dare la definizione di valori regolari e autovalori di $ A $, si consideri la seguente
quantità, sempre definita:
$$ C := \inf\limits_{\ket{\psi}\in\mathcal{D}_A}\frac{||(A-a\mathbb{I})\ket{\psi}||}{||\ket{\psi}||}\geq 0$$ 
\textbf{Def.} $a\in\mathbb{R}$ si dice punto regolare di $ A $ se $ C > 0 $.

\textbf{Def.} $a\in\mathbb{R}$ si dice autovalore di $ A $ se $C=0$

Poiché $ C $ è un estremo inferiore si distinguono due casi:
\begin{itemize}
    \item $ C $ è anche il minimo: in questo caso $ a $ è un \textbf{autovalore proprio} di $ A $
	ed $ \exists \ket{a}\in\mathcal{D}_A $ tale che $ A\ket{a} = a\ket{a}$, detto
	\textbf{autovettore proprio}.
    \item $ C $ non è il minimo: in questo caso $ a $ è un \textbf{autovalore generalizzato} di 
	$ A $ e l'autovettore corrispondente $ \ket{a}\notin\mathcal{D}_A $ è detto
	\textbf{autovettore generalizzato} . 
\end{itemize}

L'insieme degli	autovalori propri è detto \textbf{spettro discreto di A} \textit{Spd}(A).

L'insieme degli autovalori generalizzati è detto \textbf{spettro continuo di A} \textit{Spc}(A).

L'importanza degli spettri di un operatore è evidenziata dal seguente risultato, detto appunto
\textit{teorema spettrale}.

\noindent\fbox{%
    \parbox{\textwidth}{%
    \textbf{Teorema spettrale}
    
    Dato un operatore $ A $ autoaggiunto, le soluzioni dell'equazione $ A\ket{a} = a\ket{a} $ sono
    un insieme \textbf{completo}.

    $ \forall f \in S $ si definiscono 
\begin{itemize}
    \item I coefficienti $ f_i := \bra{a_i}\ket{f} $, con $ \ket{a_i} $ autovettori propri
	corrispondenti agli autovalori

	$ a_i \in$ \textit{Spd}(A);
    \item La funzione $ f(a):=\bra{a}\ket{f} $, con $ \ket{a} $ autovettori generalizzati 
	corrispondenti agli autovalori $ a\in $ \textit{Spc}(A).
\end{itemize}

Si può scrivere $ \forall f,g\in S $
\begin{flalign*}
 & \bra{f}\ket{g} = \sum_{i}^{}f^*_ig_i + \int_{Spd(A)}^{}daf^*(a)g(a)
    := \sumint_{Sp(A)}daf^*(a)g(a)\text{ , prodotto scalare}&&\\
& \ket{f} = \sumint_{Sp(A)}daf(a)\ket{a}\text{ , rappresentazione di $\ket{f}$}&&\\
& \sumint_{Sp(A)}\ket{a}\bra{a} = \mathbb{I} \text{ , relazione di completezza dello spettro} &&
\end{flalign*}


    }%
}

Il problema di questi risultati è che valgono solo per i vettori $\ket{f}\in S$,
ma possono essere generalizzati ricordando che $ S $ è denso in
$ \mathcal{H} = L^2(\mathbb{R}) $: questo implica che
$$ \forall\ket{\psi}\in\mathcal{H},\; \exists \left\{ \ket{\psi_n} \right\}\subset S
\text{ t.c. } \lim\limits_{n\to\infty}\ket{\psi_n} = \ket{\psi} $$
\newpage
Un'altra nozione utile è quella di \textit{limite in media quadratica}

\textbf{Def.} Limite in media quadratica di una successione di funzioni di prova
definita su $ I\subset\mathbb{R} $:
$$ \lim\limits_{n\to\infty}^{m.q.} \psi_n(x) = \psi(x) \leftrightarrow 
\lim\limits_{n\to\infty} \int_{I}^{}dx \left| \psi_n(x)-\psi(x)\right|^2 = 0$$
\textbf{Oss.} L'integrale è alla Lebesgue ed è quindi definito a meno di funzioni
quasi ovunque nulle: questo non dà però ambiguità, in quanto i prodotti scalari non
variano al variare del rappresentante della classe di equivalenza ed essi sono l'unica
parte della teoria che fa predizioni fisiche.

Date queste definizioni, si considerano ora le relazioni del teorema spettrale con
$ \ket{\psi}\in\mathcal{H}$: per lo spettro discreto non ci sono problemi, in quanto
$ \bra{a_i}\ket{\psi} $ è sempre un prodotto scalare in $ \mathcal{H} $. Per lo spettro
continuo invece $ \psi(a) = \bra{a}\ket{\psi} $ non è ben definita, perché è l'azione
della distribuzione $ \bra{a} $ su un vettore di $ \mathcal{H} $: si definisce dunque
$ \psi(a):=\lim\limits_{n\to\infty}^{m.q.}\bra{a}\ket{\psi_n} $, con
$ \left\{ \ket{\psi_n} \right\} \subset S$ successione che converge a
$ \ket{\psi}\in\mathcal{H} $.

Con queste accortezze, si può utilizzare l'intero spettro di un operatore $ A $ come base
per lo spazio di Hilbert $ \mathcal{H} $ di un sistema fisico.

\textbf{Esempio 1} L'operatore posizione in 1D $ X $.

La sua equazione agli autovalori è $ X\ket{x_0} = x_0\ket{x_0} $: proiettandola nella 
sua rappresentazione nello spazio delle coordinate si ottiene
$$ \bra{x}X\ket{x_0} = x_0\bra{x}\ket{x_0} \to x\bra{x}\ket{x_0} = x_0\bra{x}\ket{x_0} $$
$ x_0 $ è un parametro, quindi $ \bra{x}\ket{x_0} := f(x) $; nessuna funzione di $ L^2 $
soddisfa $ xf(x) = x_0f(x) $, ma la $\delta(x-x_0)$ di Dirac, come limite in media quadratica,
sì.

Dunque gli autovettori di $X$ sono generalizzati e la loro rappresentazione nello spazio delle
coordinate è $ \bra{x_0}\ket{x} = \delta(x-x_0) $,
inoltre lo spettro è il continuo $ x_0\in\mathbb{R} $. 

Vale la relazione di completezza $ 1=\int_{-\infty}^{\infty}dx\ket{x}\bra{x} $ e dunque
i prodotti scalari possono essere rappresentati come
$ \bra{f}\ket{g} = \int_{-\infty}^{\infty}dx \bra{f}\ket{x}\bra{x}\ket{g} =
\int_{-\infty}^{\infty}dxf^*(x)g(x)$ 

\textbf{Esempio 2} L'operatore numero d'onda in 1D $ K $ 

L'operatore $ K $ è quell'operatore che ha come rappresentazione nello spazio delle coordinate
$ \bra{x}K\ket{x'} = \delta(x'-x)\left( -i\dv{x} \right) $; la sua equazione agli autovalori è
$$ K\ket{k} = k\ket{k} \to \bra{x}K\ket{k} = k\bra{x}\ket{k} \to
\int_{-\infty}^{\infty}dx'\bra{x}K\ket{x'}\bra{x'}\ket{k} = k\bra{x}\ket{k} $$
che per la definizione di $ K $ e identificando $ \bra{x}\ket{k} := f_k(x) $ diventa
$$ -i\dv{f_k}{x} = kf_k(x) \leftrightarrow f_k(x) = N e^{ikx}$$
$ f_k(x) $ è un autovettore generalizzato e lo spettro è continuo ($ k\in\mathbb{R} $): 
$ f_k(x) $ può essere normalizzato a
$ f_k(x) = \frac{1}{\sqrt{2\pi}}e^{ikx} $, in modo che $ \bra{k}\ket{k'} = \delta(k-k') $. 

Per il teorema spettrale, $ \ket{f} = \int_{-\infty}^{\infty}dk\bra{k}\ket{f}\ket{k}
= \int_{-\infty}^{\infty}dkf(k)\ket{k}$ e si
può ricavare, ricordando la relazione di completezza, che

$$ f(x)=\bra{x}\ket{f} = \int_{-\infty}^{\infty}dk\bra{x}\ket{k}\bra{k}\ket{f} = 
\frac{1}{\sqrt{2\pi}}\int_{-\infty}^{\infty}dk e^{-ikx}f(k)$$ 


\subsection{Sistemi Completi di Osservabili Commutanti}
Uno dei concetti più importanti per risolvere problemi di meccanica quantistica è
quello di Sistema Completo di Osservabili Commutanti (SCOC), cioè un insieme di operatori
$ \{ A_1,...,A_n \} $ tutti commutanti tra di loro e tali che
$\forall B\text{ t.c. } [B,A_i] = 0,\: B =f(A_i,A_2,...,A_n)$: il massimo
numero di operatori commutanti $ n $  è il numero di gradi di libertà del sistema quantistico.

Detti $ A_i $ gli operatori del SCOC di un sistema quantistico,$\exists\left\{\ket{\alpha}\right\}$ 
set di autostati comuni a tutti gli $ A_i$:
ciascun vettore $ \ket{\alpha} $ può essere univocamente identificato dai suoi $ n $
autovalori $ \alpha_i $, con $ A_i\ket{\alpha} = \alpha_i\ket{\alpha} $, come conseguenza
del seguente teorema.

\noindent\fbox{%
    \parbox{\textwidth}{%
    \textbf{Teorema}, gli autospazi simultanei del SCOC hanno dimensione 1.

    \textbf{Dimostrazione} Si assuma, per assurdo, che $ \ket{v_1}, \ket{v_2} $
    (base dell'autospazio degenere) abbiano gli stessi autovalori $ \alpha_i $ per gli operatori 
    $A_i$ e si definisca l'operatore $B$ indipendente dal SCOC t.c. $ B\ket{v_1} = \ket{v_2}$
    e $B\ket{v_2} = \ket{v_1}$:
     allora $ A_iB\ket{v_1} = A_i\ket{v_2} = \alpha_i \ket{v_2}$ e
     $ BA_i\ket{v_1} = \alpha_i\ket{v_2} $ e similmente con $ \ket{v_2} $; questo significa
     che nell'autospazio degenere $[B,A_i] = 0$, dunque che $ B = F(A_1,...,A_n)$ e che
     $ B\ket{v_1} = f(\alpha_1,...,\alpha_n)\ket{v_1} = \ket{v_2}$, che è assurdo,
     in quanto $ \ket{v_1},\ket{v_2}$ sono linearmente indipendenti. 

     Dunque $ \nexists \ket{v_1},\ket{v_2}$ linearmente indipendenti con stessi $\alpha_i$,
     cioè non c'è degenerazione. \textbf{QED}

    }%
}

Da qui in avanti si denoterà $ \ket{\alpha} = \ket{\alpha_1,\alpha_2,...,\alpha_n}$.

\textbf{Oss.} Poiché il sistema è completo,
$\sumint d\alpha_1...d\alpha_n \ket{\alpha_1,...,\alpha_n}\bra{\alpha_1,...,\alpha_n} =\mathbb{I}$

\subsection{Esempio noto: la buca di potenziale finita 1D}
Come riepilogo delle nozioni formali appena fornite, si considera ora il noto esempio della
buca rettangolare finita, definita dal potenziale
$$
V(x) = \begin{dcases}
 0   & \text{se } |x| > a \\
 V_0   & \text{se } |x| < a
\end{dcases}
$$
con $ a>0 $ e $ V_0 < 0 $. Poiché $ V(x) $ è indipendente dal tempo, considero l'equazione
di Schrödinger stazionaria
$$ H\ket{\psi} = E\ket{\psi} \to \frac{-\hbar^2}{2m}\dv[2]{\psi}{x} + V(x)\psi(x)= E\psi(x)
    \to \dv[2]{\psi}{x} + \frac{2m}{\hbar^2}(E-V(x))\psi = 0
$$ 
Per notazione, si definiscono I, II, III gli intervalli $ (-\infty,-a],\,[-a,a],\,[a,\infty)$. 

La risoluzione dell'equazione differenziale si divide in 3 casi, a seconda del valore di $E$:

\begin{itemize}[leftmargin=*]
\item $ \mathbf{E>0}$

Definiti $ k_I^2 = \frac{2mE}{\hbar^2} $ e $ k_{II}^2 = \frac{2m(E-V_0)}{\hbar^2} $, si ha
che l'equazione differenziale è

$$
\begin{dcases}
    \psi'' + k_I^2\psi = 0 &|x| >a \\
    \psi'' + k_{II}^2\psi = 0 & |x|<a
\end{dcases}
$$
 che porta alla soluzione

 $$ 
\psi(x) = \begin{dcases}
    Ae^{ik_Ix} + Be^{-ik_Ix} &x<-a\\
    Ce^{ik_{II}x} + De^{-ik_{II}x} &x\in[-a,a]\\
    Fe^{ik_Ix} + Ge^{-ik_Ix} &x>a
\end{dcases}
 $$
 con $ A,B,C,D,F,G\in\mathbb{C} $ costanti. Imponendo la continuità di $ \psi $ e $ \psi' $ in
 $ -a $ e $ a $, si trovano 4 equazioni lineari con 6 incognite: il sistema è sempre
 risolubile a meno di due parametri liberi e indipendenti dal valore di $ E,\,V_0 $.

 Questo evidenzia che per $E>0$  lo spettro di $ H $ è continuo, con $ E\in\mathbb{R}^+ $,
 ma anche che è presente una degenerazione di grado 2 (per i due parametri liberi):
 deve dunque esistere un operatore che commuti con $ H $, e questo è l'operatore parità
 $ \Pi $, che ha come autospazi gli insiemi delle funzioni pari o dispari;
 questa distinzione tra gli autovettori di $ H$ è evidente scrivendo gli esponenziali complessi
 come somme di seni e coseni
 $$ Ae^{ik_Ix} + Be^{-ik_Ix} = a_1\cos(k_Ix) + a_2\sin(k_Ix) \quad x<-a$$
 che possono produrre funzioni a parità definita; il SCOC è dunque $ \{H,\Pi\} $ e gli stati
 del sistema possono essere identificati dal \textit{ket} $ \ket{E,\pm} $, dove $ \pm $ indica
 la parità dell'autofunzione.

 Questo risultato è generale, per il seguente teorema

 \noindent\fbox{%
     \parbox{\textwidth}{%
	 \textbf{Teoremino}
	 Dati due stati $\ket{\psi_1},\ket{\psi_2}$ tali che $H\ket{\psi_i} = E\ket{\psi_i}$, con
	 $[H,\Pi]=0$, si possono sempre costruire autostati di $H$ simmetrici o antisimmetrici.

	 \textbf{Dim.} Considero $\ket{\pm}:=\ket{\psi_i}\pm\Pi\ket{\psi_i}$:

	 $\Pi\ket{\pm}= \Pi\ket{\psi_i}\pm\Pi^2\ket{\psi_i} = \Pi\ket{\psi_i}\pm\ket{\psi_i} =\pm\ket{\pm}$,
	 cioè $\ket{\pm}$ è, rispettivamente, simmetrico o antisimmetrico; inoltre
	 $H\ket{\pm}= H\ket{\psi_i}\pm H\Pi\ket{\psi_i} = E(\ket{\psi_i} \pm \Pi\ket{\psi_i}) = E\ket{\pm}$
	 \textbf{QED}
     }%
 }
 
 Un altro modo comune di risolvere la degenerazione è indicare se la particella incida sulla buca
 da destra o da sinistra, ponendo, rispettivamente, $ A=0 $ o $ G=0 $.
 %itemize ends in next file

  
